\subsection{Modern Approaches}\label{subsec:modern}

In addition to the classical local thresholding methods derived from Niblack's formulation~\cite{Niblack1985}, we also investigate two more recent adaptive binarization approaches that emphasize computational efficiency and robustness under challenging illumination conditions.

% Integral Image: Fast Mean-Based Approach
\paragraph{Integral-Image-Based Thresholding.}
Bradley and Roth~\cite{Bradley2007} propose a fast adaptive thresholding technique based on integral images.
The central idea is to compute local mean pixel values (intensities) efficiently by using a cumulative sum representation.

% Integral Processing
Given an image $f(x,y)$, the integral image $S(x,y)$ is defined as 

\begin{equation}
  S(x,y) = \sum_{i=0}^{x} \sum_{j=0}^{y} f(i,j).
\end{equation}

Using this representation, the sum of intensities within any rectangular window can be obtained in constant time, independent of the window size.
This significantly reduces the computational cost of sliding-window thresholding.

The method compares each pixel value against the average intensity of its surrounding $s \times s$ neighborhood.
A pixel is classified as foreground if its intensity is at least $t$ percent lower than the local mean.
Otherwise, it is assigned to the background.
The overall procedure consists of the following steps:

\begin{enumerate}
  \item Compute the integral image representation of the input.
  \item Apply the adaptive mean-based threshold to produce the binarized output.
\end{enumerate}

% Bataineh: Naive Local Statistics Approach
\paragraph{Adaptive Statistics Thresholding.}
Bataineh et al.~\cite{Bataineh2011} introduce an adaptive thresholding approach that incorporates both local and global image statistics. 
In contrast to purely mean-based methods, their formulation adjusts the threshold using an adaptive standard deviation term derived from local window variations.

The threshold is defined as

\begin{equation}
  T = m_W - \frac{m^2_W - \sigma_W}{(m_g + \sigma_W) \times (\sigma_{\text{Adaptive}} + \sigma_W)} ,
\end{equation}

where $m_W$ and $\sigma_W$ denote the mean and standard deviation within the current window, and $m_g$ is the global mean intensity of the entire image.

The adaptive standard deviation is computed as a normalized measure of local contrast:

\begin{equation}
  \sigma_{\text{Adaptive}} = \frac{\sigma_W - \sigma_{\min}}{\sigma_{\max} - \sigma_{\min}},
\end{equation}

where $\sigma_{\min}$ and $\sigma_{\max}$ represent the minimum and maximum window standard deviations observed across the image.
This normalization allows the method to adjust the threshold dynamically depending on the degree of local variation.
